% ============================================================================
% CAPÍTULO 5 — RESULTADOS
% Objetivo: presentar lo que la tesis efectivamente demuestra.
% Resultado que DEBE contener (esto ES el núcleo defendible):
%   · Trayectorias de colapso R(tau) validadas en el régimen de polvo.
%   · El mapa de condiciones iniciales (retrato de fase + potencial).
%   · Clasificación de regímenes y cruce del horizonte.
%   · La condición de black shell (a=1/2) ilustrada sobre el mapa.
%   · (Si se hizo) resultados del régimen a>1 con el paso adaptativo.
% Regla: cada figura/afirmación acá debe trazarse a un objetivo del cap. 1.
%   Si una figura no sirve a un objetivo, no va en resultados.
% ============================================================================
\chapter{Resultados}
\label{cap:resultados}

\section{Colapso en el régimen de polvo}
% TODO.

\section{Mapa de condiciones iniciales}
% TODO: leer la Fig.~\ref{fig:mapa}; qué muestra que una trayectoria sola no.

\section{Cruce del horizonte y condición de black shell}
% TODO.
